\section{Experiments}
\label{app:experiments}

This appendix contains additional details on experiments from the main paper, in addition to
other experiments ran to supplement the findings in this work. For the conditional reverse task,
we take the opportunity to highlight insights we omitted from the main paper in light of space constraints.
For the other experiments, we give brief summaries of the findings.

The code for analyses and generating the figures are available in the supplementary material. They are
a useful reference for finer details (Appendix \ref{app:implementation}).

\subsection{Conditional Reverse}
\label{app:experiments-cond-rev}

\subsubsection{Models}

The transformers we use (similar to Fig. \ref{fig:conditional_reverse}) have their MLP layer before attention. This is atypical. But for the conditional reverse
task it is the more natural ordering. The MLP can tag tokens with information about whether they are vowels,
and attention can leverage that information to perform the conditional reverse. We tried the typical ordering
but \t{Model I} struggles to solve the task when limited to a one-layer transformer.

Subsequent experiments (e.g., rotate and conditional rotate) test transformers with a natural ordering,
and with two layers of transformers.

\subsubsection{Does compiling change learning?}

In Fig. \ref{fig:behavioral-dynamics}, the heatmaps for each model show a horizontal line toward the midpoint of the y-axis.
Each model learns the correct output token at the 8th position faster than the rest. This is because for the conditional
reverse task on a sequence of 15 characters, the 8th position is always identity. The task at that position is simpler.

\subsubsection{Does compiling change representations?}

In Fig. \ref{fig:pca}, each PCA plot shows representations projected onto the first 3 principal components (PCs). In general,
these three components tend to explain around 50\% of variance in each plot. PC1 is the bottom-left axis. PC2 is the bottom-right axis.
PC3 is the vertical axis. The exact amount of variance explained by each is below:

\begin{table}[ht]
\centering
\begin{tabular}{llccc}
\toprule
\textbf{Model} & \textbf{Representation} & \textbf{PC1} & \textbf{PC2} & \textbf{PC3} \\
\midrule
F & R2 & 45.4\% & 19.7\% & 14.1\% \\
        & A  & 47.9\% & 23.3\% & 10.1\% \\
\midrule
P & R2 & 21.1\% & 16.8\% & 13.7\% \\
        & A  & 20.6\% & 15.3\% & 14.1\% \\
\midrule
T & R2 & 25.4\% & 20.4\% & 18.6\% \\
        & A  & 22.9\% & 19.6\% & 18.1\% \\
\midrule
I & R2 & 18.3\% & 16.2\% & 14.4\% \\
        & A  & 18.4\% & 16.2\% & 14.4\% \\
\bottomrule
\end{tabular}
\end{table}

\subsubsection{Does compiling change the necessary representations?}

One observation we did not mention about Fig. \ref{app:experiments-cond-rev} is that \t{Model T} seems to use both the compiled and non-compiled 
attention heads. Recall that its compiled head is the $\rd{k}$-linear map denoted by an encoding of the
conditional reverse program in Fig. \ref{fig:conditional_reverse}---but where the coefficients of that map
are randomized (e.g., think randomly initialized matrix or tensor). It does not relearn the solution
denoted by \t{Model F}, even though there exists a configuration of its coefficients which implement \t{Model F}.

\subsubsection{Does compiling change the controllable representations?}

An interesting detail about our encoding of characters in \t{Model F} is that it is \textit{not} a
one-hot encoding. We encode characters using the type $\nb{\+_{i\bl{=}0}^{9}}\Unit = \nb{\Unit \+ \Unit \+ \cdots \+ \Unit}$. This type
can exactly implement an alphabet of 10 characters using one-hot encodings. But for an alphabet of
26 characters, our encoding uses the first coordinate to represent whether a character is a vowel,
and the other 9 coordinates are essentially left to the neural network to decide how to use. An important
consequence of our steering results in Fig. \ref{fig:steering} is that transformers can learn to respect
this encoding.


\subsection{Reverse}
\label{app:experiments-rev}

The findings of this experiment are consistent with what is observed in the main paper. There are two
main distinctions worth noting. First, \t{Model F} is now compiled from a reverse program, essentially \t{Model P} 
from the main paper. Moreover, there is a different ablation \t{Model U}, which is essentially \t{Model F} but
where we \textit{unfreeze} the coefficients of its $\rd{k}$-linear map. We were interested in whether gradient descent would decide to preserve
the compiled circuit, or opt to rewire it into something else. For this experiment, it appeared to preserve
the compiled circuit. But in subsequent experiments (e.g., rotate and conditional rotate) \t{Model U} exhibited
extremely unstable learning dynamics. 


\begin{figure}[h]
  \centering

  % Left half: one large plot
  \begin{subfigure}[c]{0.48\linewidth}
    \centering
    \includegraphics[width=\linewidth]{../experiments/reverse/data/plots/behavioral_dynamics_lr0.01_bs128_maxstep200.pdf}
    \vspace{.3em}
    \caption{\normalsize Behavioral Dynamics}
  \end{subfigure}
  \hfill
  % Right half: two stacked plots
  \begin{minipage}[c]{0.48\linewidth}
    \centering

    \begin{subfigure}{\linewidth}
      \centering
      \begin{tabular}{c}
        \t{qahftrxckafnafq|qfanfakcxrtfhaq}\\
        \t{ofpvausieyiccwp|pwcciyeisuavpfo}\\
        \t{usnzjovqwpsbfhc|chfbspwqvojznsu}\\
        \t{gchqjjfgyqpesej|jesepqygfjjqhcg}
      \end{tabular}
    \vspace{.3em}
      \caption{\normalsize Task}
    \end{subfigure}

    \vspace{2.5em}

    \begin{subfigure}{\linewidth}
      \centering
      \includegraphics[width=\linewidth]{../experiments/reverse/data/plots/ablations_step2000_lr0.01_bs128.pdf}
        \vspace{.3em}
      \caption{\normalsize Resample Ablations}
    \end{subfigure}

  \end{minipage}
\end{figure}

\begin{figure}[h] 
  \centering 
  \includegraphics[width=.85\linewidth]{../experiments/reverse/data/plots/pca_res_step2000_lr0.01_bs128_seed0.pdf}
  \vspace{.5em}
  \caption*{(d) PCA on Representations}
  \vspace{-1em}
\end{figure}



\begin{table}[h]
\vspace{1.5em}
\centering
\begin{tabular}{llccc}
\toprule
\textbf{Model} & \textbf{Representation} & \textbf{PC1} & \textbf{PC2} & \textbf{PC3} \\
\midrule
F & R2 & 46.8\% & 35.5\% & 8.6\% \\
  & A & 55.4\% & 22.5\% & 6.2\% \\
\midrule
U & R2 & 29.0\% & 27.5\% & 22.2\% \\
  & A & 56.7\% & 13.0\% & 10.5\% \\
\midrule
T & R2 & 42.0\% & 38.3\% & 15.0\% \\
  & A & 35.0\% & 32.0\% & 15.9\% \\
\midrule
I & R2 & 26.8\% & 21.5\% & 20.8\% \\
  & A & 26.9\% & 21.6\% & 21.0\% \\
\bottomrule
\end{tabular}
\vspace{1.5em}
\caption*{(e) Variance explained by PCs}
\end{table}

\vspace{3em}
\subsection{Rotate}
\label{app:experiments-rotate}

This experiment concerns an entirely new task: rotate. The transformers must learn to shift the input
sequence to the left $n$ times, where $n$ is the position of the first character in the alphabet (0-indexed).
For example, a string starting with \t{a} is just identity because \t{a} is in the $0$-th position in the
alphabet. A string starting with \t{b} moves to the left once (e.g., \t{bac} becomes \t{acb}). This task
is more complicated than reverse, so each model is a two-layer transformer with the standard ordering, attention
layers coming before MLP layers. The compiled attention head encoding rotate is only present in the first layer of the transformer.

Because \t{Model T} and \t{Model U} do not learn the task, our analyses focus on \t{Model F} and \t{Model I}.

\begin{figure}[h]
  \centering

  % Left half: one large plot
  \begin{subfigure}[c]{0.48\linewidth}
    \centering
    \includegraphics[width=\linewidth]{../experiments/rotate/data/plots/behavioral_dynamics_lr0.001_bs128_maxstep5000.pdf}
    \vspace{.3em}
    \caption{\normalsize Behavioral Dynamics}
  \end{subfigure}
  \hfill
  % Right half: two stacked plots
  \begin{minipage}[c]{0.48\linewidth}
    \centering

    \begin{subfigure}{\linewidth}
      \centering
      \begin{tabular}{c}
        \t{qahftrxckafnafq|ahftrxckafnafqq}\\
        \t{ofpvausieyiccwp|pofpvausieyiccw}\\
        \t{usnzjovqwpsbfhc|ovqwpsbfhcusnzj}
      \end{tabular}
    \vspace{.3em}
      \caption{\normalsize Task}
    \end{subfigure}

    \vspace{2.5em}

    \begin{subfigure}{\linewidth}
      \centering
      \includegraphics[width=\linewidth]{../experiments/rotate/data/plots/ablations_step5000_lr0.001_bs128.pdf}
    \vspace{.3em}
      \caption{\normalsize Resample Ablations}
    \end{subfigure}

  \end{minipage}
\end{figure}

\begin{figure}[h] 
  \centering 
  \includegraphics[width=.55\linewidth]{../experiments/rotate/data/plots/pca_res_step5000_lr0.001_bs128_seed0.pdf}
  \vspace{.8em}
  \caption*{(d) PCA on Representations}
  \vspace{-1em}
\end{figure} 



\begin{table}[h]
\vspace{1.5em}
\centering
\begin{tabular}{llccc}
\toprule
\textbf{Model} & \textbf{Representation} & \textbf{PC1} & \textbf{PC2} & \textbf{PC3} \\
\midrule
F & R1 & 12.7\% & 10.8\% & 10.0\% \\
  & A & 30.9\% & 20.8\% & 16.5\% \\
\midrule
I & R1 & 22.6\% & 12.7\% & 9.5\% \\
  & A & 38.6\% & 21.4\% & 15.3\% \\
\bottomrule
\end{tabular}
\vspace{1em}
\caption*{(e) Variance explained by PCs}
\end{table}




\clearpage
\subsection{Conditional Rotate}
\label{app:experiments-cond-rotate}

The final experiment is the most difficult. It is a variant of the previous rotate task. If the last character is a vowel,
you rotate left, and otherwise rotate right by $n$ positions
where $n$ is the position of the first character in the alphabet (0-indexed). Model architectures are similar to the rotate
experiment. However, this time a program encoding the conditional rotation is compiled into the attention head of the first
transformer layer.

Because \t{Model T} and \t{Model U} do not learn the task, our analyses focus on \t{Model F} and \t{Model I}.

\begin{figure}[h]
  \centering

  % Left half: one large plot
  \begin{subfigure}[c]{0.48\linewidth}
    \centering
    \includegraphics[width=\linewidth]{../experiments/cond_rotate/data/plots/behavioral_dynamics_lr0.001_bs128_maxstep5000.pdf}
    \vspace{.3em}
    \caption{\normalsize Behavioral Dynamics}
  \end{subfigure}
  \hfill
  % Right half: two stacked plots
  \begin{minipage}[c]{0.48\linewidth}
    \centering

    \begin{subfigure}{\linewidth}
      \centering
      \begin{tabular}{c}
        \t{qahftrxckafnafq|qqahftrxckafnaf}\\
        \t{ofpvausieyiccwp|fpvausieyiccwpo}\\
        \t{usnzjovqwpsbfhc|sbfhcusnzjovqwp}
      \end{tabular}
    \vspace{.3em}
      \caption{\normalsize Task}
    \end{subfigure}

    \vspace{2.5em}

    \begin{subfigure}{\linewidth}
      \centering
      \includegraphics[width=\linewidth]{../experiments/cond_rotate/data/plots/ablations_step5000_lr0.001_bs128.pdf}
    \vspace{.3em}
      \caption{\normalsize Resample Ablations}
    \end{subfigure}

  \end{minipage}
\end{figure}

\begin{figure}[h] 
  \centering 
  \includegraphics[width=.45\linewidth]{../experiments/cond_rotate/data/plots/pca_res_step5000_lr0.001_bs128_seed0.pdf}
  \vspace{.8em}
  \caption*{(d) PCA on Representations}
  \vspace{-1em}
\end{figure} 

\begin{table}[h]
\vspace{1.5em}
\centering
\begin{tabular}{llccc}
\toprule
\textbf{Model} & \textbf{Representation} & \textbf{PC1} & \textbf{PC2} & \textbf{PC3} \\
\midrule
F & R1 & 14.8\% & 12.6\% & 11.1\% \\
  & A & 39.2\% & 15.6\% & 11.9\% \\
\midrule
I & R1 & 14.8\% & 10.2\% & 9.2\% \\
  & A & 32.8\% & 23.4\% & 14.1\% \\
\bottomrule
\end{tabular}
\vspace{1em}
\caption*{(e) Variance explained by PCs}
\end{table}